%%==========================================================
%%  CAPÍTULO 4: DISEÑO Y CONSTRUCCIÓN DEL MECANISMO DE GENERACIÓN DE PROYECTOS
%%==========================================================
\label{cap:generacion}

El Capítulo~\ref{cap:fundamentacion} encuadró la solución como línea de
producto de software y fijó los fundamentos sobre los que se apoya. Un
encuadre, sin embargo, no reduce por sí solo el esfuerzo de iniciar un
proyecto; lo reduce un mecanismo que lo aplique sin intervención manual. Esta
fase construyó ese mecanismo, y con él las reglas concretas que el marco dejó
asignadas, que son qué forma tiene un proyecto válido de la familia y por qué
contratos se le acopla una capacidad compartida.

Lo que entregó es una herramienta de línea de comandos instalable que, a
partir de los activos de la línea y de una selección de características,
produce un proyecto backend conforme al marco. La condición de llegada se
formuló antes de construir y es verificable, porque exige un producto derivado
que arranca sin que nadie edite un archivo, que supera su propio análisis
estático con la regla de dependencia activa como error, y que declara con qué
versión de cada activo fue producido. El detalle de la construcción, con sus
decisiones intermedias y la evidencia de cada verificación, se conserva en los
documentos de la Fase 2, disponibles en el repositorio del proyecto.

\section{Hitos de la construcción}
\label{sec:hitos-fase-2}

La fase se ejecutó en cinco hitos, y el orden entre ellos es en sí mismo una
decisión, porque los que construyen los activos van antes que el que construye
el generador. Todo lo que un activo contiene puede verificarse a mano mientras
no exista todavía ningún generador, de modo que invertir el orden habría
obligado a descubrir un error de la estructura cuando ya existía una
herramienta edificada alrededor de ella. La Tabla~\ref{tab:hitos-generacion}
recoge qué dejó cada hito y cómo quedó comprobado.

\begin{table}[H]
\centering
{\footnotesize
\begin{tabular}{p{2.9cm}p{6.1cm}p{5.4cm}}
\toprule
\textbf{Hito} & \textbf{Qué dejó construido} & \textbf{Cómo quedó comprobado} \\
\midrule
La forma del producto & Las capas, la regla de dependencia, la frontera del contrato externo y los contratos de integración de la línea & El análisis estático falla nombrando la regla cuando se introduce una importación prohibida \\
La plantilla base & Un proyecto ejecutable con el entorno de trabajo configurado y la documentación de API automática & Arranca sin que nadie edite un archivo y supera su propia verificación \\
El catálogo & El manifiesto declarativo, el modelo de gobierno y la distribución por el registro de la organización & Se publica desde su repositorio y se consume en ejecución, sin plataforma propia \\
La herramienta & El núcleo común con sus seis puertos, el punto de extensión y la operación de derivación & La suite corre sin tocar disco, red ni terminal, con dobles de cada puerto \\
El almacén de plantillas & Una plantilla que se describe a sí misma, el motor de gramática cerrada y la forma del producto como contrato & Una segunda plantilla, de otro modelo de ejecución, entró sin que cambiara una línea de la herramienta \\
\bottomrule
\end{tabular}
}
\caption{Hitos de la construcción del mecanismo de generación}
\label{tab:hitos-generacion}
\end{table}

El quinto hito no estaba previsto al abrir la fase, y se abrió al comprobar que
el alcance comprometía un almacén de plantillas, en plural, y que la línea
tenía una sola plantilla sin lugar donde poner una segunda; la
sección~\ref{sec:almacen} narra qué obligó a revisar y qué quedó más firme
que antes.

\subsection{Alcance de la fase}
\label{sec:alcance-fase-2}

La frontera con la fase siguiente no es temporal sino funcional, y esa
distinción evita el riesgo principal del diseño, porque ambas fases construyen
la misma herramienta. Lo que toda operación comparte se construye aquí, y lo que
es propio de una operación se construye en la fase que la introduce. En
particular, el punto de extensión por el que se agregan operaciones nuevas
tiene que existir antes de que la Fase 3 agregue las suyas; de lo contrario
esa fase rehace la herramienta en lugar de ampliarla. La
Tabla~\ref{tab:alcance-fase-2} recoge lo que queda fuera y por qué.

\begin{table}[H]
\centering
{\footnotesize
\begin{tabular}{p{6.2cm}p{8.8cm}}
\toprule
\textbf{Queda fuera} & \textbf{Razón} \\
\midrule
Incorporación de módulos a un producto & El marco trata generar y componer como operaciones distintas sobre el mismo mecanismo \\
Construcción de módulos del catálogo & La composición se valida en la fase donde el módulo tiene función \\
Actualización sobre productos existentes & Requiere más de una versión publicada de un activo consumido \\
Despliegue e infraestructura & Excluidos en la especificación de requerimientos \\
\bottomrule
\end{tabular}
}
\caption{Alcance diferido de la fase}
\label{tab:alcance-fase-2}
\end{table}

Todo lo diferido se registra como punto de extensión y no como limitación de
la solución, que es la diferencia entre una frontera decidida y una carencia.

\section{La forma de un producto de la familia}
\label{sec:organizacion-codigo}

Los fundamentos adoptados establecen qué contiene el núcleo de un producto y
cómo se relaciona con el exterior, pero no dónde reside físicamente cada
elemento. Eso lo fijó la línea antes que ninguna otra cosa, porque sin una
forma declarada no hay nada que derivar.

El proyecto se organiza primero por capa y, dentro de la capa de dominio, por
agregado.

\subsection{Las reglas de dependencia, y por qué son verificables}

Las dependencias entre capas admiten una sola dirección, porque la
infraestructura puede depender de la aplicación y del dominio; la aplicación,
únicamente del dominio; y el dominio, de ninguna de las dos. La restricción no
alcanza a dos elementos, que son el punto de composición, donde se asocian los
puertos con sus adaptadores, y el contrato de la línea, que solo declara
tipos.

De ahí la regla que la línea adopta como criterio: ningún elemento de dominio
importa un elemento de aplicación o de infraestructura. El dominio no conoce
el framework, el motor de persistencia ni el protocolo de transporte.

Lo que convierte esa regla en una restricción efectiva no es enunciarla sino
cómo se comprueba. Se verifica por análisis estático sobre las rutas de
importación, no por revisión manual, y la violación se reporta como error y
no como advertencia, de modo que un proyecto que la incumpla no supera su
verificación. Una regla que solo advierte no restringe.

Las convenciones de nomenclatura que acompañan a esta estructura, con el
sufijo que corresponde a cada patrón táctico y la forma de derivar un nombre
desde el estándar, se recogen en los documentos de la Fase 2.

\subsection{La frontera del contrato externo y la documentación de API}

Los tipos que un adaptador de entrada expone al exterior son distintos de los
objetos de dominio, y el adaptador traduce explícitamente entre unos y otros,
porque son dos contratos con ciclos de vida independientes. El objeto de
dominio cambia cuando cambian las reglas del negocio; el tipo expuesto, solo
cuando cambia el acuerdo con quienes consumen el servicio. Si el adaptador
expusiera el objeto de dominio, cualquier reorganización interna se
convertiría en un cambio incompatible para los consumidores.

Esa separación tiene una consecuencia aprovechable que cierra un
requerimiento, porque los tipos expuestos, al ser explícitos, son la fuente
desde la cual se deriva la documentación de API del producto, porque el
documento se arma con los esquemas que el producto ya usa para validar sus
peticiones y respuestas, y no con una declaración paralela. Es lo que impide
que la documentación y el comportamiento se separen, porque no hay dos fuentes
que mantener sincronizadas sino una. El documento lleva además, para cada
operación, el identificador y el grupo que el estándar le da, que es lo único
que permite leerlo contra el publicado cuando la forma de las rutas del
producto puede variar.

La interfaz de exploración se decidió dos veces, y la segunda decisión no
contradijo las razones de la primera. Al construir el documento se descartó un
visor del tipo \emph{Swagger UI}, porque exigiría una dependencia de ejecución
en todo producto o un script traído de internet, y ninguna de las dos es
admisible en el entorno al que la línea sirve. Lo que esa decisión dejaba a
cada equipo, la línea lo resolvió después una sola vez, de modo que el
producto sirve una ruta que renderiza el mismo documento con el visor cargado
de la copia instalada con el producto y nunca de internet, con la dependencia
declarada como de desarrollo. Así no cambia lo que un servicio desplegado
carga, la auditoría de dependencias de ejecución sigue dando el mismo
resultado, y una variable de entorno decide si la ruta existe, apagada por
defecto en producción.

\section{Puertos y contratos de integración}
\label{sec:puertos-contratos}

La estructura admite que la capa de aplicación dependa de contratos de
puerto, pero no dice qué son ni quién es su dueño. Resolverlo es la condición
de todo el mecanismo de composición, porque de ahí depende cómo un activo
reutilizable se acopla a un producto sin romper el desacoplamiento que la
regla de dependencia establece. Es también el punto donde la arquitectura
hexagonal se sostiene o se vuelve nominal, porque si el módulo declarara el
puerto, el producto acabaría dependiendo del módulo y la dirección quedaría
invertida en la práctica, aunque la definición afirmara lo contrario.

\subsection{El contract package}

Un módulo debe poder servir a más de un producto, y para ello alguien tiene
que declarar el puerto. La solución es un artefacto neutral del que dependen
tanto los productos como los módulos, y del que ninguno de los dos depende
del otro. Es el \emph{contract package}, un \emph{core asset} que contiene
las declaraciones de los puertos técnicos de la línea.

Lo que hace viable esta decisión es que la dependencia de un producto hacia él
es exclusivamente de tipos, y una importación de tipos se elimina durante la
compilación, de modo que el producto declara el contrato pero en ejecución no
existe referencia alguna a él. Sin esa propiedad sería una dependencia real y
la solución perdería la mitad de su valor.

De ahí su regla de contenido, que es la ausencia de todo aquello que
sobreviva a la compilación. No se enuncia como «solo interfaces», dado que
una interfaz necesita tipos para sus parámetros y retornos. Cuando un puerto
opera sobre datos del dominio de un producto se declara genérico y el
producto aporta el tipo, con lo que el contrato nunca conoce el dominio de
ningún producto y el producto nunca depende del módulo.

\subsection{Las dos clases de puerto}

El criterio que separa las dos clases de la Tabla~\ref{tab:puertos} es la
firma y no el uso: un puerto es de dominio si su firma menciona conceptos del
dominio del producto, y es técnico si su firma es expresable sin conocer ese
dominio.

\begin{table}[H]
\centering
{\footnotesize
\begin{tabular}{p{2.4cm}p{4.8cm}p{3.5cm}p{4.0cm}}
\toprule
\textbf{Clase} & \textbf{Criterio} & \textbf{Quién lo declara} & \textbf{Quién lo realiza} \\
\midrule
De dominio & Su firma menciona conceptos del dominio del producto, como el repositorio de un agregado & El producto, en su capa de dominio & El propio producto, en su infraestructura \\
Técnico & Su firma es expresable sin conocer el dominio, como el registro de eventos o el cifrado & El \emph{contract package} & Un módulo del catálogo \\
\bottomrule
\end{tabular}
}
\caption{Las dos clases de puerto y la propiedad de cada una}
\label{tab:puertos}
\end{table}

La consecuencia es que ningún módulo puede apropiarse de un puerto de
dominio. Un repositorio expresa intención de negocio en el lenguaje del
producto, y no puede provenir de un activo genérico sin arrastrar el dominio
hacia el catálogo. En sentido inverso, un módulo sí puede proponer un puerto
técnico nuevo, porque exigir aprobación previa para cada uno convertiría el
contrato en el cuello de botella del catálogo. El caso que obliga a ejercitar
las dos clases a la vez es la persistencia, y se desarrolla en el
Capítulo~\ref{cap:gestor}, donde el módulo que la realiza se construye.

\section{La variabilidad de la línea}
\label{sec:variabilidad}

El modelo de variabilidad declara qué puede diferir entre un producto y otro.
Es el contrato que la herramienta consume, y por eso se escribe en un formato
declarativo legible por máquina, porque sin él la parametrización se decidiría
sobre la marcha durante la implementación y dejaría de ser verificable.
Reside en artefactos versionados, cada uno con su esquema, que la
verificación de la línea valida en cada entrega.

Las decisiones se agrupan en las tres clases de la
Tabla~\ref{tab:clases-variabilidad}, y la clase determina la forma de la
respuesta.

\begin{table}[H]
\centering
\begin{tabular}{p{2.6cm}p{5.4cm}p{6.6cm}}
\toprule
\textbf{Clase} & \textbf{Qué resuelve} & \textbf{Forma de la decisión} \\
\midrule
Identidad & Qué representa el producto y cómo se nombra & Un valor escrito o elegido \\
Selección & Qué capacidades del dominio incorpora & Un subconjunto de un conjunto declarado \\
Presencia & Qué piezas opcionales se incluyen & Inclusión u omisión \\
\bottomrule
\end{tabular}
\caption{Clases de variabilidad del modelo}
\label{tab:clases-variabilidad}
\end{table}

La clase de selección es la que sitúa este modelo dentro del cuadrante de
análisis de requerimientos de la aplicación, porque es el punto donde el
desarrollador elige características y no solo nombra lo que se genera. El
modelo declara trece puntos de variación, cada uno con su clase, su valor por
defecto y las condiciones que lo habilitan o lo excluyen.

\subsection{El espacio de selección de dominio}

El \emph{Service Domain} no se escribe libremente, sino que se elige de un
conjunto declarado. La escritura libre dejaría el metadato de dominio sin
validación posible y convertiría la alineación con el estándar en una
afirmación no verificable, que es precisamente lo que el marco buscaba
evitar. Cada entrada describe una capacidad de negocio con su área, su grupo,
la entidad de información que gobierna y las operaciones que el estándar
publica sobre ella. El estándar modela cada capacidad en dos niveles y el
espacio recoge ambos, de modo que al derivar se ofrecen tanto las operaciones
que gobiernan el ciclo de vida de la capacidad como las de sus subdivisiones.

Una entrada existe cuando sirve a un propósito escrito del proyecto, y se
agrega junto con ese propósito. No se agrega para poblar un catálogo ni para
cubrir un \emph{Business Domain}, porque transcribir entero el subconjunto
que el marco declaró convertiría una lista de referencia en una asignación
definitiva y pondría a elegir entre capacidades que el contexto no va a
construir.

\paragraph*{La fuente del dato, y su límite}

El dato de una capacidad se transcribe del documento que el estándar publica
para ese \emph{Service Domain}, y no de una vista general del
\emph{Service Landscape} ni de una paráfrasis. La distinción no es formal, y lo
demostró una transcripción hecha sin ese documento, que atribuyó a una
capacidad una cláusula que pertenece a otra, le asignó una operación en un
nivel que el estándar no usa y declaró como no mutante una operación que el
documento
publica con cuerpo de petición. Los tres errores se corrigieron al conseguir
la fuente.

La consecuencia es un límite declarado y no un hueco oculto, porque solo las
capacidades cuyo documento publicado está a la vista se transcriben, y las
demás conservan su descripción parafraseada hasta que lo esté. Inventar el
contenido del estándar sería el modo de perder exactamente la propiedad que
justifica adoptarlo.

\paragraph*{La capacidad que el catálogo no ofrece}

El mecanismo admite declarar un \emph{Service Domain} fuera del espacio por
defecto, registrando el \emph{Business Domain} que lo contiene, y prohíbe
únicamente omitir el metadato. El contexto lo necesita porque no todo backend
del laboratorio corresponde a una capacidad que el estándar publique.

Lo que se mantiene cerrado es la \emph{forma} de la respuesta, no siempre el
conjunto de valores. Una capacidad declarada elige su \emph{Business Domain}
entre los que el marco reconoce y sus operaciones entre los términos de
acción que el mecanismo deriva, y el \emph{Business Area} se lee del grupo en
lugar de preguntarse, lo que impide que un producto responda las dos cosas
por separado y se contradiga con el estándar. El metadato sigue validado, de
modo que lo que cae es la forma y no la razón.

El registro del producto marca el caso, y esa marca habilita una pregunta que
antes no tenía dónde formularse: cuántos productos de la línea responden por
una capacidad que el estándar no publica, y cuáles de ellas se repiten. Tres
productos declarando la misma capacidad es la evidencia que justifica
incorporarla al catálogo, y constituye el flujo de retorno de requerimientos
hacia el dominio que el modelo de referencia declara.

\section{De la plantilla base al almacén de plantillas}
\label{sec:almacen}

\subsection{La plantilla base como validación del marco}
\label{sec:activo-base}

La plantilla base es el proyecto del que se derivan todos los demás, y el
primer artefacto ejecutable de este trabajo. Materializa la estructura de la
sección~\ref{sec:organizacion-codigo} y trae el entorno de trabajo ya
configurado, con el análisis estático y su regla de dependencia activada como
error, el ejecutor de pruebas, las herramientas de formato y la configuración
externalizada que RNF-05 exige.

Su función va más allá de servir de punto de partida, porque es la validación
del marco. Si no arranca, o no supera su propio análisis estático, el problema
está en la definición y no en el código, y hay que corregir la definición
antes de seguir. Esa es la razón por la que se construye antes que la
herramienta.

Una decisión de este hito condiciona el resto de la fase, y es que la plantilla
base no se mantiene a mano, sino que es exactamente el resultado de rendir las
plantillas con una entrada determinada. La verificación lo comprueba archivo
por archivo, con la única excepción de los tres que llevan la identidad del
artefacto, que en el activo describen un activo y en un producto describen un
producto. La consecuencia es que no pueden divergir, porque un cambio en una
plantilla que no se refleje en el activo rompe la verificación y un retoque
manual del activo la rompe igual. Es el mecanismo que impide que la plantilla
y su referencia se separen en silencio, que es el modo habitual en que un
generador deja de producir lo que su ejemplo muestra.

\subsection{Una plantilla es un activo que se describe a sí mismo}

Con una sola plantilla, el reparto entre lo que la plantilla declara y lo que
la herramienta sabe no tiene que decidirse, porque todo puede vivir en la
herramienta sin que se note. El alcance del proyecto compromete otra cosa, ya
que habla de un sistema de almacenamiento de plantillas y RF-14 del catálogo
de plantillas y módulos, en plural. Hacer exigible esa promesa obligó a hacer
el reparto explícito, y lo que salió es más firme que lo que había.

Una plantilla pasa a ser un activo del catálogo del mismo rango que un
módulo, y como tal se describe a sí misma en un manifiesto que viaja en su
paquete, declarando el contrato contra el que se construyó, qué forma de
producto deriva, dónde viven sus archivos y de qué está hecho su producto. El
manifiesto va en el paquete y no en el catálogo por el mismo argumento que
rige para un módulo, porque lo que una plantilla declara de sí misma cambia
con ella, y ponerlo en el catálogo obligaría a coordinar dos repositorios para
publicar una versión. Sus archivos salieron de la herramienta al paquete de
la plantilla, con lo que una plantilla de otro equipo puede vivir en el
repositorio de ese equipo y una versión de la plantilla deja de costar una
versión de la herramienta.

\paragraph*{La forma del producto, como contrato}

La línea había descrito en prosa los puntos donde un módulo inserta su
cableado, y la herramienta los conocía como constantes de su propio código.
Con una plantilla eso bastaba; con dos deja de bastar, porque una plantilla
que ponga su punto de composición en otro archivo sería inservible. La forma
del producto pasó entonces a ser un dato con nombre y versión, declarado en
tres lugares que se comparan entre sí: la plantilla declara qué forma derivan
sus productos, el producto la lleva escrita en su registro, y el módulo
declara a qué formas sirve.

Un punto de inserción de los ocho que la forma hexagonal declara no dice solo
dónde se inserta, sino con qué sintaxis se comenta el archivo que lo contiene,
que es lo que permite que el mismo motor escriba en un archivo de TypeScript y
en uno de configuración sin conocer ninguno de los dos.

El producto lleva la forma y no se lee de la plantilla, por la misma razón
por la que registra con qué versión del activo fue derivado. La
compatibilidad de un módulo se decide contra el producto, y la plantilla que
lo derivó puede no estar instalada, no estar publicada o no existir. Lo que
el producto es tiene que poder leerse del producto.

El cambio tuvo además una consecuencia de la clase que una definición en prosa
no obliga a resolver. El marco suponía que
el producto había derivado del activo que escribió esas marcas, y nadie lo
comprobaba. Hoy son dos rechazos de la composición: uno cuando un módulo
reemplaza una región que la forma del producto no declara, y otro cuando al
producto le falta un punto que su propia forma declara. Un supuesto implícito
pasó a ser comprobable por los dos lados.

\subsection{El motor y la frontera de lo expresable}

Con un solo activo base el motor podía limitarse a sustituir nombres, y la
propiedad que lo protegía era que las plantillas no deciden nada. Un almacén
con más de una plantilla obliga a revisar esa propiedad, porque una plantilla
que no pueda expresar una condición tiene que duplicarse en dos archivos por
cada decisión, y eso fue lo que llevó a que una sola plantilla ocupara más de
cincuenta archivos.

El motor ganó entonces condición y repetición, pero con una gramática
cerrada, que es una propiedad más fuerte que la anterior y no una renuncia.
Antes se afirmaba que las plantillas no deciden nada; ahora se afirma que
todo lo que una plantilla puede decidir está enumerado, en tres
construcciones y dos tablas cerradas. No hay evaluación dinámica de código,
no existe forma de que una plantilla nombre una función, y la condición la
lee el mismo evaluador que ya leía las restricciones del modelo de
variabilidad, de modo que no hay una segunda gramática que mantener. Una
etiqueta que el motor no reconoce es un defecto de la plantilla y aparece en
la verificación del activo que la posee. Con esa capacidad, los archivos de
plantilla se redujeron casi a la mitad.

La frontera de lo expresable se escribió como parte de la decisión, y dos
cosas se eliminaron en lugar de expresarse. El formato por ancho de línea
salió porque es formato y lo comprueba el verificador del propio producto, y
el impresor de literales del documento de API pasó a ser texto en su
plantilla. Lo que queda del lado de la herramienta queda por una razón
declarada, porque es mecanismo y no producto: las huellas de integridad, el
reloj, el respaldo, las preguntas, el registro, el catálogo y los códigos de
salida.

El efecto sobre el código fue que la función que escribía el producto dejó de
ser el catálogo de archivos del producto escrito en código y pasó a ser un
intérprete, con una firma que no nombra ni el framework, ni el estándar, ni
ninguna carpeta de un producto.

\subsection{La evidencia de que el mecanismo es genérico}

Todo lo anterior descansa en una proposición: que lo que queda en la
herramienta es un mecanismo y no la descripción de un ejemplo. Con una sola
plantilla la proposición no es comprobable, porque un intérprete genérico y
una función escrita para ese caso se comportan igual.

La evidencia se produjo en dos formas que prueban cosas distintas, y la primera
es una plantilla de otra familia como fixture de prueba, en otro lenguaje,
con otro estilo de comentario, un trabajo por lotes en lugar de un servicio y
un contrato que no es el de esta línea; prueba que el intérprete deriva una
familia de entrada distinta. La segunda es una plantilla publicable que
deriva un API HTTP sin servidor propio, con otra forma de producto; prueba lo
que el fixture no puede, que es que una plantilla de otra forma pasa las
puertas de la línea, deriva un producto que registra esa plantilla y esa
forma, se instala y supera su propia verificación.

Lo que la segunda plantilla difiere no es un segundo caso de lo mismo. No
tiene contenedor de inyección, no tiene servidor propio, y sus puntos de
inserción son cinco en lugar de ocho, con uno de los roles que sencillamente
no existe. Ese último dato es el hallazgo del contrato de forma, porque el
archivo de configuración de esa plataforma es JSON estricto y un punto de
inserción es una línea de comentario que ahí dejaría el archivo inválido. Una
forma cuyo archivo de configuración es JSON no puede tener punto de entorno, y
el contrato lo expresa sin inventar nada, porque la forma declara los puntos
que tiene y el mecanismo rechaza nombrando el que falta.

La ronda dejó además probado lo que era una condición de diseño declarada y
nunca ejercida, que es que el diseño no quedó atado a un framework concreto.
Una plantilla sobre otro modelo de ejecución entró sin que una línea de la
herramienta cambiara.

\section{Gobierno de los activos y catálogo}
\label{sec:gobierno-activos}

Una línea de producto sostiene tres actividades en paralelo: el desarrollo de
los activos reutilizables, el de los productos que se derivan de ellos, y la
gestión que decide el alcance, asigna responsables y sostiene la línea en el
tiempo \citep{Clements2001}. Las dos primeras quedaron resueltas con los
activos y el mecanismo; la tercera es la que responde a la brecha B-03, porque
sin gestión los activos derivan, nadie responde por ellos y el catálogo se
convierte en una lista de cosas que alguien publicó alguna vez.

\subsection{Tipos de activo y escala de adaptación}

El catálogo reconoce cuatro tipos de \emph{core asset}, de los que la
plantilla aporta la estructura de partida de un producto; el módulo, la
realización de una o más capacidades técnicas; el \emph{contract package}, las
declaraciones de los puertos técnicos de la línea; y el manifiesto del
catálogo describe a los demás. Tres condiciones son comunes a todos: un
responsable declarado, un versionado que distinga el cambio incompatible del
compatible, y la posibilidad de saber qué productos lo consumen. La tercera es
la que hace gobernable al conjunto, porque permite responder qué se ve
afectado por un cambio consultando el catálogo, sin abrir los repositorios de
los productos. El responsable es un equipo y no una persona, porque un activo
cuyo responsable fuera una persona quedaría huérfano al salir esa persona de
la organización. Las condiciones específicas por tipo las detalla el documento
de gobierno de los activos de la Fase 2.

El equipo que recibe el producto derivado conserva la autonomía de
modificarlo y, al mismo tiempo, la línea debe poder propagarle cambios. La
escala de adaptación de la Tabla~\ref{tab:adaptacion} resuelve esa tensión
declarando tres niveles con consecuencias distintas.

\begin{table}[H]
\centering
{\footnotesize
\begin{tabular}{p{2.8cm}p{5.6cm}p{6.3cm}}
\toprule
\textbf{Nivel} & \textbf{Qué hace el equipo} & \textbf{Efecto sobre la propagación} \\
\midrule
Configuración & Ajusta valores previstos por el activo, fuera del código & El producto sigue recibiendo toda mejora del activo \\
Reemplazo & Sustituye el adaptador de un puerto por otro propio, conservando el contrato & El producto sigue recibiendo las mejoras que no tocan el adaptador reemplazado \\
Copia desvinculada & Copia el código del activo al producto y lo modifica & El producto deja de recibir mejoras, y el catálogo lo registra como desvinculado \\
\bottomrule
\end{tabular}
}
\caption{Escala de adaptación y su efecto sobre la propagación de cambios}
\label{tab:adaptacion}
\end{table}

Los dos primeros niveles convierten la adaptabilidad en una propiedad
arquitectónica y no en una concesión. El tercero se admite en lugar de
prohibirse por una razón práctica, porque prohibirlo no evitaría que ocurra,
solo lo dejaría sin registro. Una copia desvinculada se lee entonces como
señal de que al activo le falta un punto de configuración o de reemplazo, y
alimenta el flujo de nuevos requerimientos hacia el análisis del dominio.

\subsection{Versionado y cambio incompatible}

La línea fija qué constituye un cambio mayor, uno menor y un parche, porque
de ello depende que dos productos puedan operar con versiones distintas de un
mismo activo, que es el criterio de RNF-06. El documento de gobierno de los
activos de la Fase 2 enumera cada caso.

Lo que importa es el criterio con que se decide, porque la compatibilidad se
determina por inspección de lo declarado, es decir la firma de las
operaciones, el esquema de configuración y el identificador de registro. Toda
la evidencia necesaria está en el contrato, sin recurrir a la ejecución, y
con el lenguaje adoptado el compilador detecta por sí mismo los cambios de
firma que rompen a un consumidor. De ello se sigue una limitación asumida de
forma explícita, y es que el activo no se comprueba ejecutándolo, de modo que
un cambio de comportamiento que no altere la firma pasa inadvertido. Se registra
como punto de extensión.

El \emph{contract package} se organiza en submódulos por capacidad pero se
versiona como una unidad, de modo que un cambio incompatible en una capacidad
eleva la versión mayor de todas. La consecuencia se asume, porque un producto
que consume una sola capacidad recibe una versión mayor que no le afecta, y
un módulo puede verse obligado a republicarse sin que su implementación haya
cambiado. La separación en artefactos independientes queda prevista, con una
señal de cuándo conviene hacerla: cuando un módulo deba republicarse sin
alterar su implementación más de dos veces, el acoplamiento de versión cuesta
más de lo que ahorra en gobernanza.

\subsection{El catálogo, su gobernanza y la resolución en ejecución}
\label{sec:catalogo-operacion}

El catálogo es un manifiesto declarativo y versionado que reside en un
repositorio propio y que \emph{describe} activos sin contenerlos: una entrada
nombra el paquete y su responsable, de modo que un módulo construido por otro
equipo vive en el repositorio de ese equipo y el catálogo solo lo registra.
Esa forma es lo que satisface RF-14 sin desplegar ni operar ninguna
plataforma, y lo que cierra la brecha B-03 al alcance de un equipo pequeño.
Los campos de una entrada y los estados por los que pasa se recogen en ese
mismo documento.

Incorporar una entrada exige que el activo satisfaga las condiciones de
validez y que la propuesta sea revisada por un rol responsable del gobierno
de la línea, cuya existencia define el marco mientras que a quién se le
asigna lo decide cada organización que lo adopte. No se admite retirar un
activo, porque uno del que dependen productos existentes no puede eliminarse
sin romperlos. El estado de desuso es terminal, el activo deja de ofrecerse al
derivar y permanece en el catálogo indefinidamente. El retiro acompañado
de un procedimiento de migración se registra como trabajo futuro.

La forma en que la herramienta obtiene el catálogo cambió al medir su costo, y
el episodio ilustra un criterio. La primera versión fijaba la versión del
catálogo dentro de la herramienta, con el argumento de que RF-05 exige que la
misma entrada produzca el mismo resultado. Medido, ese argumento no se
sostenía y el costo era alto, porque una entrada nueva costaba dos
publicaciones y una de ellas era una versión de la herramienta cuyo único
cambio sustantivo era un número. El determinismo que RF-05 pide no lo daba la
versión del catálogo sino el pin de la plantilla, porque lo que determina la
estructura de un producto es el contenido del activo del que se deriva, no la
lista de activos que existen. El catálogo pasa entonces a resolverse en
ejecución, con tres fuentes en orden —la ruta que dé la invocación, la última
versión que sirve el registro, y la copia instalada cuando el registro no
responde—, informando cuál de ellas respondió.

Lo que sí queda fijado dentro de la herramienta es la versión exacta de cada
plantilla que esa entrega ofrece, con una versión exacta y nunca un rango,
porque un rango haría que dos ejecuciones de la misma herramienta resolvieran
versiones distintas. Un módulo no se fija así, sino que se resuelve a la
última publicada y la versión obtenida queda registrada en el producto, ya
que pinar cada módulo obligaría a publicar una versión de la herramienta por
cada versión de cada módulo, que es exactamente la dependencia que el
catálogo existe para evitar.

\section{El núcleo de la herramienta y la operación de derivación}
\label{sec:herramienta-cli}

Con los activos construidos y el catálogo en operación, el hito siguiente es
el mecanismo que los consume. La herramienta se construyó separando lo que
toda operación comparte de lo que es propio de una, y esa separación es lo
que permitió a la fase siguiente ampliarla en lugar de rehacerla.

\subsection{El núcleo y sus puertos}
\label{sec:nucleo-herramienta}

Pertenece al núcleo aquello que una segunda operación necesitaría igual: la
resolución de la entrada contra el modelo de variabilidad, la consulta al
catálogo, la lectura y escritura del registro de versiones del proyecto, la
conversión de toda escritura en un plan de cambios antes de aplicarlo, y la
traducción de cualquier fallo a un código de salida estable.

El marco exige arquitectura hexagonal al proyecto generado, y la herramienta
la aplica a sí misma. Su núcleo no conoce el sistema de archivos, el registro
de paquetes, la terminal ni el reloj, sino que los declara como los seis
puertos de la Tabla~\ref{tab:puertos-nucleo} y recibe sus implementaciones
desde fuera.

\begin{table}[H]
\centering
\begin{tabular}{p{4.8cm}p{10.2cm}}
\toprule
\textbf{Puerto} & \textbf{Qué abstrae} \\
\midrule
Sistema de archivos & Leer, escribir y comprobar existencia de rutas \\
Registro de paquetes & Obtener un activo y resolver la dirección del registro de un ámbito \\
Preguntas & Pedir al desarrollador lo que la invocación no trajo \\
Reporte & Lo que la herramienta dice, con su forma \\
Reloj & El momento que el producto registra como su procedencia \\
Ejecutor de comandos & Invocar procesos externos, como el gestor de paquetes o el control de versiones \\
\bottomrule
\end{tabular}
\caption{Los seis puertos del núcleo de la herramienta}
\label{tab:puertos-nucleo}
\end{table}

La consecuencia práctica es que el comportamiento del núcleo se prueba sin
tocar disco, sin red y sin terminal, con dobles escritos a mano que
implementan cada puerto. Esa decisión es la que hace que la suite corra en
segundos y que una prueba falle por el motivo que dice y no por el entorno
donde se ejecuta.

Ese contrato es estrecho, porque una operación declara su nombre, su resumen,
los parámetros que admite y una función que recibe la entrada ya validada
junto con el contexto de puertos.
El despachador no sabe qué operaciones existen, sino que las lee de un
registro, de modo que agregar una es escribir su módulo y registrarla sin
modificar el núcleo. Ese punto de extensión es el entregable que la fase
siguiente consume, y su existencia se comprobó antes de necesitarla, porque la
operación de consulta del catálogo se construyó en esta misma fase sobre el
mismo contrato, de modo que la segunda operación existió mientras el punto de
extensión todavía podía corregirse.

\subsection{La operación de derivación}
\label{sec:comando-generacion}

La operación de derivación resuelve la plantilla contra el catálogo y la
obtiene del registro, recoge las decisiones de variabilidad, comprueba que la
combinación elegida satisfaga las reglas del modelo, y convierte todo ello en
un plan de cambios que el núcleo escribe. Al terminar existe un proyecto que
arranca sin que nadie edite un archivo.

Una operación de negocio elegida no produce un archivo sino un recorrido
completo desde la superficie externa hasta el contrato del repositorio: un
método en el adaptador HTTP con su ruta, un caso de uso en la capa de
aplicación, el método que ese caso de uso necesita en el puerto del
repositorio, y un evento de dominio cuando la operación cambia el estado de la
entidad que la capacidad gobierna. Ese reparto no se decide archivo por
archivo, sino que cada operación recorre el mismo camino desde que se elige
hasta los archivos que produce, pasando por la entrada que la herramienta
resuelve y por las reglas que convierten un término de acción en un nombre de
clase, una ruta, un método y un evento. Es lo que hace que la selección de
características tenga consecuencias estructurales y no cosméticas.

\paragraph*{El límite de lo que se genera}

El mecanismo produce la estructura, los nombres y los contratos vacíos,
mientras que las reglas de negocio las construye el equipo que recibe el
proyecto, y mantener esa distinción es lo que mantiene al artefacto dentro de
su alcance.

La distinción se declara donde se nota, de modo que las operaciones cuyo
efecto se sigue del propio término de acción se generan completas; las que
dependen de la capacidad se generan cargando la entidad, persistiéndola,
publicando su evento y dejando el efecto marcado con un comentario. El
registro del producto declara, por operación, cuál de los dos casos es, y el
producto lo dice en su propia documentación. Un producto que pareciera más
terminado de lo que está sería el peor resultado posible de un generador.

\paragraph*{Determinismo}

La misma entrada produce el mismo producto, sin importar quién ejecute el
comando. Tres decisiones lo sostienen: la versión de la plantilla se resuelve
contra el pin que viaja en la herramienta y no contra el registro; toda
escritura pasa por un plan de cambios que es un dato inspeccionable antes de
aplicarse; y el único valor que varía entre dos ejecuciones idénticas, que es
el momento de la derivación, se obtiene por un puerto y se registra en un
campo conocido. La verificación deriva el mismo producto dos veces y compara
los árboles resultantes excluyendo ese campo.

\paragraph*{El recorrido de entrada}

La operación cierra con una revisión que muestra cada decisión tomada y
ofrece corregir cualquiera de ellas antes de escribir nada. Corregir una
respuesta vuelve a preguntar únicamente lo que depende de ella, y la razón es de
adopción y no de comodidad, porque una pregunta del estándar contestada mal
costaba el comando entero, y ese costo es el que hace que un desarrollador
deje de leer las preguntas. Fuera de una terminal la revisión declara y
continúa, porque una corrida automatizada no tiene a quién preguntar.

El recorrido completo se corrigió además en un punto que ninguna prueba veía,
y la causa es estructural, porque los circuitos automáticos corren con
banderas y ningún camino automático pasa por una pregunta. Lo que apareció al
recorrerlo a mano fue que una condición no satisfecha se comprobaba después
de cerrar todas las preguntas, y el costo de no satisfacerla era perderlas
todas. Hoy la condición viaja con la pregunta y se rechaza ahí mismo, la
elección de plantilla ocurre antes de la primera pregunta del modelo, y la
tecla de escape retrocede una pregunta en lugar de abandonar el comando.

El recorrido abre eligiendo la plantilla, antes de cualquier pregunta del
modelo, ofreciendo las dos opciones del almacén con lo que cada una deriva
como pista y mostrando la lista entera en lugar de enrollarse sobre sí misma.

\section{Verificación y conformidad}
\label{sec:pruebas-fase-2}

El último hito no agrega funcionalidad, sino que produce la evidencia de que
lo construido hace lo que debía, y opera en dos niveles cuya distinción es
deliberada. El primero prueba la herramienta con dobles de sus puertos,
comprobando la resolución de la entrada, la evaluación de las reglas del
modelo, el reparto de archivos por operación y los nombres derivados, sin
tocar disco ni red. El segundo deriva productos reales y los somete a su
propia verificación —análisis estático, pruebas y compilación—, que es la
única forma de saber que lo generado es código válido y no texto bien
formado.

\subsection{Romper la regla para saber que existe}

Un producto que supera su análisis estático demuestra que el analizador
corre, no que la regla de dependencia esté activa. La verificación introduce
por eso, a propósito, una importación que la regla prohíbe, exige que el
análisis falle nombrando esa regla, y restaura el archivo. El criterio se
generalizó a toda regla nueva de la línea: una regla cuya violación no se
haya comprobado no está verificada, está escrita.

El criterio dio su propio resultado, porque una regla que exigía una prueba de
humo por cada ruta declarada se satisfacía con que la ruta apareciera en el
título de la prueba, no con que alguna petición la ejercitara, de modo que una
prueba apuntada a otra ruta la habría pasado. El defecto apareció al intentar
romperla y no al escribirla.

\subsection{Conformidad del producto derivado}

El producto se verifica contra la condición de llegada de la fase y no contra
una lista de archivos: arranca sin edición previa, responde sus sondas,
registra lo que sirve, no registra las sondas, y no contiene direcciones ni
lecturas del entorno fuera de su módulo de configuración. El circuito deriva
varios productos con entradas distintas —capacidades distintas, las dos
formas de ruta, con y sin eventos, una capacidad declarada fuera del
catálogo— y cada uno pasa su propia verificación, que es la evidencia que el
Capítulo~\ref{cap:evaluacion} cita.

La documentación de instalación y uso de la herramienta vive en el
repositorio de los activos y viaja en el paquete publicado, no en este
documento, por una razón de mantenimiento, y es que una guía de uso transcrita
en la tesis quedaría desactualizada en la primera versión publicada.

\section{Resumen del capítulo}

Esta fase convirtió un encuadre en un mecanismo, de modo que lo que al cerrar el
Capítulo~\ref{cap:fundamentacion} era una definición es ahora una herramienta
instalable que produce proyectos conformes a ella sin intervención manual, y
con ella las tres reglas concretas que el marco había dejado asignadas: la
forma de un producto válido, los contratos por los que se le acopla una
capacidad compartida, y el gobierno de los activos que lo sostiene en el
tiempo.

Cuatro decisiones estructurales sostienen el resultado, y la primera es el
orden de los hitos, porque los activos se construyeron antes que el
generador, de modo que la plantilla base funcionó como validación ejecutable
del marco cuando corregirlo todavía era barato. La segunda es la separación
entre el núcleo común y la operación concreta, con un punto de extensión
comprobado antes de necesitarlo, que es lo que permite a la fase siguiente
ampliar la herramienta en lugar de rehacerla. La tercera es que la plantilla
base es la salida de sus propias plantillas y no un artefacto mantenido en
paralelo, lo que hace imposible que se separen en silencio. Y la cuarta es el
\emph{contract package} como artefacto neutral del que dependen productos y
módulos sin depender entre sí, que es la pieza de la que depende todo el
mecanismo de composición del capítulo siguiente.

La variabilidad quedó declarada en artefactos versionados y no en el código
del generador, y esa separación es la que permitió después ampliar el espacio
de selección, incorporar el segundo nivel del estándar y admitir una
capacidad fuera del catálogo sin modificar la herramienta. Es también el
punto donde el estándar interviene en el mecanismo y no solo en el registro,
porque sin el espacio de selección la alineación con BIAN se reduciría a un
metadato escrito a mano.

El quinto hito dejó la lección más incómoda y la más útil, porque con un solo
activo base, el reparto entre lo que la plantilla declara y lo que la
herramienta sabe podía no decidirse, y la consecuencia fue que la herramienta
guardaba el catálogo de archivos del producto dentro de su código. Cerrar la
promesa de un almacén obligó a hacer ese reparto explícito, y lo que salió es
más firme que lo que había: una plantilla que se describe a sí misma, un motor
cuya capacidad de decisión está enumerada, una derivación que es un intérprete
y una forma de producto comprobable por los dos lados. La evidencia de que el
resultado es genérico no es un argumento sino una segunda plantilla, de otro
modelo de ejecución, que entró sin que una línea de la herramienta cambiara.

Queda abierto, por construcción, todo lo que exige módulos que aún no
existen: incorporarlos a un producto, consultarlos, actualizarlos y preservar
el código propio del equipo al hacerlo. Es el objeto del
Capítulo~\ref{cap:gestor}.
